\documentclass[11pt]{article}
\usepackage[margin=1in]{geometry}
\usepackage{amsmath,amssymb}
\usepackage{booktabs}
\usepackage{longtable}
\usepackage{array}
\usepackage{graphicx}
\usepackage{hyperref}
\usepackage{enumitem}
\usepackage{xcolor}
\usepackage{float}
\usepackage{placeins}
\usepackage{caption}

\captionsetup{font=small, labelfont=bf, width=0.95\textwidth}

\title{A Sparse Poisson GLM-HMM for Simultaneous Neural Connectivity\\and Brain-State Recovery\\[0.5em]
\normalsize Full Mathematical Description and Implementation Reference}
\author{Sanjit Rao}
\date{April 2026}

\begin{document}
\maketitle

% ============================================================
\section*{Abstract}
% ============================================================

A central problem in systems neuroscience is recovering the functional connectivity
of a neural circuit --- which neurons drive which --- from observed spike trains alone.
This is hard for two reasons: (1)~neural dynamics are \emph{non-stationary},
with brain states (arousal, attention, task phase) modulating baseline excitability
and masking connectivity signals; and (2)~indirect multi-hop pathways produce
lag-1 correlations that are difficult to distinguish from direct connections.

The \textbf{GLM-HMM model} described here addresses both challenges simultaneously.
It models spike counts as a \textbf{Poisson GLM-HMM}: a generative model in which
each time bin is assigned to one of $K$ discrete brain states, and within each state
the spike rate is a log-linear function of the previous bin's population activity via
a shared connectivity matrix $A$.  Three structural constraints derived from biology
are enforced throughout training: \textbf{Dale's law} (each neuron is either
purely excitatory or purely inhibitory), \textbf{spectral stability}
(spectral radius $\rho(A) < 1$), and \textbf{self-inhibition}
($A_{ii} \le -0.05$ for all neurons).

The model has been evaluated on two simulated 100-neuron networks with $T=200{,}000$
time bins ($\Delta t = 0.01$\,s).  On $K=2$ states it achieves $F_1 = 0.995$
(TP=1338, FP=13, FN=1) and state accuracy $= 0.992$.  On $K=3$ states it achieves
$F_1 = 0.997$ (TP=1358, FP=9, FN=0) and state accuracy $= 0.988$.

% ============================================================
\section*{Summary}
% ============================================================

The GLM-HMM model recovers two quantities simultaneously from neural spike-count recordings:
\begin{enumerate}[leftmargin=*]
  \item A shared recurrent connectivity matrix $A\in\mathbb{R}^{N\times N}$, where entry
        $A_{ij}$ encodes the influence of neuron $j$'s past activity on neuron $i$'s
        current firing rate.  Rows $0,\ldots,n_{\mathrm{exc}}-1$ are excitatory
        (off-diagonal entries $\ge 0$); rows $n_{\mathrm{exc}},\ldots,N-1$ are
        inhibitory (off-diagonal entries $\le 0$) --- Dale's law.
  \item A discrete hidden state sequence $z_t\in\{0,\ldots,K-1\}$ capturing non-stationary
        modes (brain state, arousal, task phase) that modulate the baseline firing rate
        of every neuron independently of the shared connectivity.
\end{enumerate}

Training uses the \textbf{Expectation-Maximisation (EM)} algorithm with \textbf{Adam}
optimisation and \textbf{gradient clipping} at norm 1.0.  Each run jointly learns
the sparse connectivity structure and the HMM parameters over 25 EM iterations.

Performance summary across evaluated datasets:

\begin{center}
\begin{tabular}{lcc}
\toprule
Metric & $K=2$ & $K=3$ \\
\midrule
F1 (overall)         & 0.995 & 0.997 \\
TP / FP / FN         & 1338 / 13 / 1 & 1358 / 9 / 0 \\
Excitatory F1        & 0.995 & 0.997 \\
Inhibitory F1        & 0.995 & 0.994 \\
Weight $R$ / slope   & 0.990 / 1.000 & 0.988 / 0.956 \\
State accuracy       & 0.992 & 0.988 \\
Bias $R$ (per state) & $\ge 0.999$ & $\ge 0.999$ \\
\bottomrule
\end{tabular}
\end{center}

Both datasets achieve near-perfect connectivity recovery ($F_1 > 0.99$).
$K=2$ achieves slope $= 1.000$; $K=3$ has slightly lower slope ($0.956$) but zero
false negatives.

% ============================================================
\section{Data and Preprocessing}
% ============================================================

\subsection{Observed arrays}
\begin{equation}
  Y\in\mathbb{R}_{\ge 0}^{T\times N},\qquad N=100.
\end{equation}
$Y_{t,n}$ is the spike count of neuron $n$ in time bin $t$.  Bins are
$\Delta t = 0.01$\,s wide.  $K$ is the number of hidden brain states; $K=2$ or
$K=3$ depending on the dataset.  Both datasets use $T = 200{,}000$ bins.

\subsection{Lagged design matrix}
\begin{equation}
  X_t =
  \begin{cases}
    \mathbf{0}, & t = 1,\\
    Y_{t-1},    & t \ge 2.
  \end{cases}
\end{equation}
$X_t\in\mathbb{R}^N$ contains spike counts from the immediately preceding bin.
Current runs use $L=1$ lag; with $L>1$ lags the design matrix is the horizontal
concatenation $X_t = [Y_{t-1},\,Y_{t-2},\,\ldots,\,Y_{t-L}]$ and $A$ grows to
shape $N\times LN$.

% ============================================================
\section{Model Parameters}
% ============================================================

\begin{equation}
  A\in\mathbb{R}^{N\times N},\quad
  B\in\mathbb{R}^{K\times N},\quad
  \pi\in\Delta^{K-1},\quad
  P\in\mathbb{R}^{K\times K}.
\end{equation}

\paragraph{Orientation.}
$A$ is stored as (post-synaptic, pre-synaptic).  With row-vector input
$X_t\in\mathbb{R}^{1\times N}$:
\begin{equation}
  \eta_{t,k} = X_t A^\top + B_k \;\in\mathbb{R}^{1\times N}.
\end{equation}
Entry $A_{ij}$ is the weight from pre-synaptic neuron $j$ to post-synaptic neuron $i$.

\paragraph{Dale's law (row convention).}
Rows $0,\ldots,n_{\mathrm{exc}}-1$ belong to excitatory neurons:
off-diagonal entries are constrained $\ge 0$.  Rows
$n_{\mathrm{exc}},\ldots,N-1$ belong to inhibitory neurons:
off-diagonal entries are constrained $\le 0$.
With $N=100$ and $\mathrm{exc\_ratio}=0.8$ this gives $n_{\mathrm{exc}}=80$.

% ============================================================
\section{Poisson Emission Model}
% ============================================================

\subsection{Clipped log-rate}
\begin{equation}
  \tilde\eta_{t,k,n} = \min(\eta_{t,k,n},\;5).
\end{equation}
The clip at 5 prevents numerical overflow in $\exp(\cdot)$ when weights or biases
are large early in training.

\subsection{Expected count}
\begin{equation}
  \mu_{t,k,n} = \exp(\tilde\eta_{t,k,n})\,\Delta t.
\end{equation}

\subsection{Emission log-likelihood}
\begin{equation}
  \ell_{t,k} = \sum_{n=1}^N
    \Bigl[Y_{t,n}\bigl(\tilde\eta_{t,k,n}+\log\Delta t\bigr) - \mu_{t,k,n}\Bigr].
\end{equation}
This is the Poisson log-likelihood up to the constant $\log(Y_{t,n}!)$.

% ============================================================
\section{Warm-Start Initialisation}
% ============================================================

\subsection{Connectivity warm start (ridge regression)}

Without a warm start, $A$ is initialised near zero and the gradient signal from the
Poisson NLL is very weak.  A ridge regression solution provides a data-informed
starting point close enough to the true $A$ that the EM optimizer can refine it
rather than search from scratch.

\begin{equation}
  A_{\mathrm{ls}} =
  \bigl(X^\top X + \lambda_{\mathrm{ridge}}\,I\bigr)^{-1} X^\top Y
  \;\in\mathbb{R}^{N_{\mathrm{in}}\times N},
  \qquad N_{\mathrm{in}} = L\cdot N.
\end{equation}
The strongest $r_{\mathrm{keep}}\cdot 100\%$ of off-diagonal entries (by magnitude) are
kept; the rest are zeroed:
\begin{equation}
  (A_{\mathrm{sparse}})_{ij} =
  (A_{\mathrm{ls}})_{ij}\,
  \mathbf{1}\!\left(|(A_{\mathrm{ls}})_{ij}|\ge
    q_{1-r_{\mathrm{keep}}}\!\left(|A_{\mathrm{ls,off}}|\right)\right),
  \quad i\neq j.
\end{equation}
Diagonal entries are set to $-0.05$ (weak self-inhibition), weights are clipped to
$[-0.5, 0.5]$, and the matrix is transposed into (post, pre) convention.
Dale's law and the inhibitory-diagonal constraint are enforced immediately after.

\subsection{Bias warm start}

If all $B_k$ are initialised identically, the forward-backward algorithm assigns equal
responsibility $\gamma_{t,k}=1/K$ to every state at every time step.  This is a saddle
point of the EM objective: the E-step produces uniform responsibilities, the M-step has
no signal to differentiate states, and the model never escapes the degenerate solution.
Small random perturbations break this symmetry.

\begin{equation}
  B_{k,n} = \log\!\left(\max\!\left(\frac{1}{T}\sum_{t=1}^T Y_{t,n},\;10^{-4}\right)\right)
            + \varepsilon_{k,n},
  \qquad
  \varepsilon_{k,n}\sim\mathcal{N}(0,\sigma_B^2),
  \quad \sigma_B = 0.05.
\end{equation}

\subsection{HMM initialisation}

Brain states are known to be persistent (dwell times of seconds to minutes), so
initialising with high self-transition probability biases the model toward
physiologically plausible state durations and reduces the risk of rapid
state-switching spurious solutions.

For $K=2$:
\begin{equation}
  \pi = [0.55,\;0.45],\qquad
  P = \begin{bmatrix}0.95 & 0.05\\0.05 & 0.95\end{bmatrix}.
\end{equation}

For $K=3$:
\begin{equation}
  \pi = [1/3,\;1/3,\;1/3],\qquad
  P = \begin{bmatrix}0.93 & 0.035 & 0.035\\0.035 & 0.93 & 0.035\\0.035 & 0.035 & 0.93\end{bmatrix}.
\end{equation}

% ============================================================
\section{E-step: Forward-Backward Algorithm}
% ============================================================

All recursions are computed in \textbf{log-space} to prevent underflow over
$T=200{,}000$ time bins.  The forward variable $\alpha_{t,k}$ is a product of $t$
probabilities each $\le 1$; log-space arithmetic replaces products with sums and
recovers numerical stability at the cost of $O(K^2)$ logsumexp operations per step.

\subsection{Forward pass}
\begin{align}
  \log\alpha_{1,k} &= \ell_{1,k} + \log\pi_k,\\
  \log\alpha_{t,k} &= \ell_{t,k}
    + \log\sum_{j=1}^K \exp\!\bigl(\log\alpha_{t-1,j}+\log P_{jk}\bigr),
    \quad t\ge 2.
\end{align}
$\alpha_{t,k}$ is the joint probability of all observations $y_1,\ldots,y_t$ and
being in state $k$ at time $t$.

\subsection{Backward pass}
\begin{align}
  \log\beta_{T,k} &= 0,\\
  \log\beta_{t,j} &= \log\sum_{k=1}^K
    \exp\!\bigl(\log P_{jk}+\ell_{t+1,k}+\log\beta_{t+1,k}\bigr),
    \quad t\le T-1.
\end{align}
$\beta_{t,k}$ is the probability of future observations given state $k$ at time $t$.

\subsection{State posteriors (responsibilities)}
\begin{equation}
  \log\gamma_{t,k} = \log\alpha_{t,k}+\log\beta_{t,k}
    - \log\sum_{j=1}^K\exp\!\bigl(\log\alpha_{t,j}+\log\beta_{t,j}\bigr),
  \qquad
  \gamma_{t,k} = \exp(\log\gamma_{t,k}).
\end{equation}
$\gamma_{t,k}$ is the posterior probability (responsibility) that state $k$ is active
at time bin $t$.

\subsection{Pairwise posteriors (for transition update)}
\begin{equation}
  \log\xi_t(i,j) = \log\alpha_{t,i}+\log P_{ij}+\ell_{t+1,j}+\log\beta_{t+1,j}
    - \log\sum_{a,b}\exp(\cdots),
\end{equation}
\begin{equation}
  \Xi_{ij} = \sum_{t=1}^{T-1}\xi_t(i,j).
\end{equation}
$\Xi_{ij}$ is the expected number of transitions from state $i$ to state $j$.

% ============================================================
\section{M-step: Sparse Structure Learning}
% ============================================================

\subsection{Loss function and optimiser}
The M-step minimises the expected negative log-likelihood under the current
responsibilities, with L2 weight decay:
\begin{equation}
  \mathcal{L} =
  \underbrace{\sum_{t\in\mathcal{B}}\sum_{k=1}^K\gamma_{t,k}
    \sum_{n=1}^N\!\bigl[\mu_{t,k,n}-Y_{t,n}(\tilde\eta_{t,k,n}+\log\Delta t)\bigr]}%
  _{\text{Poisson NLL}}
  +\;\frac{\lambda_2}{2}\|A\|_F^2.
\end{equation}
Optimisation uses \textbf{Adam} with \textbf{gradient clipping} at norm 1.0.
Gradient clipping is critical early in training: when $\rho(A)$ is large, the
recurrent dynamics amplify gradients through the emission model, causing the
optimizer to overshoot and destabilise training.  Clipping to norm 1.0 caps the
step size without discarding gradient direction information.

\subsection{Proximal soft-threshold (L1 sparsity)}

Including $\lambda_1\|A\|_1$ directly in the differentiable loss introduces a
non-zero subgradient at $A_{ij}=0$, which prevents weights from reaching exact zero
and produces a dense matrix.  The proximal operator applied \emph{after} the gradient
step produces exactly sparse solutions:
\begin{equation}
  A_{ij} \leftarrow
  \mathrm{sign}(A_{ij})\max\!\bigl(|A_{ij}|-\tau_{\mathrm{block}(i)},\;0\bigr),
  \quad i\neq j,
\end{equation}
where $\tau_{\mathrm{exc}}$ applies to excitatory rows and $\tau_{\mathrm{inh}}$ to
inhibitory rows.  Separate thresholds allow independent control of the FP/FN balance
for each block, since the two blocks have different edge densities and weight scales.

\subsection{Adaptive $\tau$ controller}

A fixed proximal threshold $\tau$ cannot adapt to the evolving weight distribution
during training.  An online feedback controller that increases $\tau$ when
$\mathrm{FP} > \mathrm{FN}$ and decreases it when $\mathrm{FN} > \mathrm{FP}$
maintains the FP/FN balance automatically without requiring per-dataset manual tuning.

After the first $25\%$ of Phase~1 epochs (warm-up period):
\begin{equation}
  e_{\mathrm{block}} = \frac{FP_{\mathrm{block}}-FN_{\mathrm{block}}}
    {\max\!\bigl(1,\;TP_{\mathrm{block}}+FP_{\mathrm{block}}+FN_{\mathrm{block}}\bigr)},
\end{equation}
\begin{equation}
  \tau_{\mathrm{block}} \leftarrow \tau_{\mathrm{block}}\cdot
    \exp\!\bigl(g_\tau\cdot e_{\mathrm{block}}\bigr),
  \qquad g_\tau = 0.05.
\end{equation}
Thresholds are clamped to
$[\tau_{\min},\tau_{\max}] = [0.2\,\tau_0,\;6.0\,\tau_0]$,
where $\tau_0=\lambda_1\cdot\alpha_{\mathrm{prox}}$ is the initial threshold.
A multiplicative (rather than additive) update is used because $\tau$ spans several
orders of magnitude across hyperparameter settings.

\subsection{Dale's law projection}

Gradient descent does not respect sign constraints.  Without projection, excitatory
neuron weights can become negative and vice versa, violating the biological constraint
that each neuron releases only one type of neurotransmitter.

\begin{align}
  A_{ij} &\leftarrow \max(A_{ij},\;0), \quad i < n_{\mathrm{exc}},\;i\neq j,\\
  A_{ij} &\leftarrow \min(A_{ij},\;0), \quad i \ge n_{\mathrm{exc}},\;i\neq j.
\end{align}

\subsection{Inhibitory diagonal constraint}

\begin{equation}
  A_{ii} \leftarrow -\max(|A_{ii}|,\;d_{\min}),\qquad d_{\min}=0.05.
\end{equation}
This ensures self-coupling is always inhibitory with minimum magnitude, preventing
runaway self-excitation and numerical instability.

\subsection{Spectral stability projection}

The recurrent dynamics are stable only if $\rho(A) < 1$.  When $\rho(A) \ge 1$,
small perturbations grow exponentially, predicted rates diverge, and the loss
becomes infinite.

\begin{equation}
  \rho(A) = \max_{\lambda\in\mathrm{eig}(A)}|\lambda|.
\end{equation}
\begin{equation}
  A[:, :N] \leftarrow A[:, :N]\cdot\frac{0.90}{\rho(A)}\quad
  \text{if }\rho(A)\ge 0.92\text{ and global\_epoch} \bmod 3 = 0.
\end{equation}

Spectral rescaling is applied every 3 epochs (not every epoch).  Applying it every
epoch creates a continuous competition between Adam (pushing weights toward true
magnitudes) and the rescaler (compressing them downward), producing permanent weight
deflation and state collapse.  Applying every 3 epochs gives Adam a window of
unconstrained gradient ascent to recover correct magnitudes before the constraint
fires again.

The rescaling target is $0.90$ (not $0.919$, the threshold).  Targeting $0.919$
leaves almost no buffer: Adam can only push $\rho$ from $0.919$ to $0.92$ before the
next rescaling fires.  A target of $0.90$ provides a buffer zone of $0.02$ in
spectral radius, allowing weights to settle closer to their true magnitudes.

A single post-training spectral stabilisation step uses a tighter target of $0.919$.

\subsection{Two-phase projection cadence}

The proximal operator and Dale's law projection are not applied every epoch.
Early in training, weights need freedom to grow before support stabilises --- applying
the proximal operator every epoch prunes edges before they have had enough gradient
signal to distinguish true edges from noise.

\begin{equation}
  q = \begin{cases} q_{\mathrm{early}} = 8 & \text{if progress} < 65\%,\\
                    q_{\mathrm{late}} = 2  & \text{otherwise.} \end{cases}
\end{equation}
Proximal and Dale projections are applied every $q$ epochs.
Spectral rescaling uses a separate fixed cadence of every 3 epochs throughout.

% ============================================================
\section{HMM Parameter Updates (M-step)}
% ============================================================

\subsection{Initial state distribution}
\begin{equation}
  \pi_k \leftarrow \frac{\gamma_{1,k}+\alpha_\pi}{\sum_j(\gamma_{1,j}+\alpha_\pi)},
  \qquad \alpha_\pi=10^{-2}.
\end{equation}

\subsection{Transition matrix}
\begin{equation}
  \tilde P_{ij} = \Xi_{ij}+\epsilon_P+\delta_{ij},
  \qquad
  P_{ij} = \frac{\tilde P_{ij}}{\sum_b\tilde P_{ib}},
  \qquad \epsilon_P=10^{-3}.
\end{equation}
The Kronecker $\delta_{ij}$ self-transition pseudocount prevents the estimated
transition matrix from collapsing to rapid state-switching when data are limited,
acting as a Bayesian prior for persistent states.

% ============================================================
\section{Final Evaluation}
% ============================================================

\subsection{Adaptive threshold search per block}
After training, the optimal binarisation threshold for each block is found by scanning
over quantiles of the non-zero inferred weights:
\begin{equation}
  \mathcal{T}_{\mathrm{block}} =
  \left\{q_p\!\left(|\hat A_{\mathrm{block}}|_{>10^{-12}}\right):
    p\in[0.70,\;0.995]\right\}.
\end{equation}
\begin{equation}
  \tau_{\mathrm{block}}^* = \arg\max_{\tau\in\mathcal{T}_{\mathrm{block}}}
    \frac{2\,TP_\tau}{2\,TP_\tau+FP_\tau+FN_\tau}.
\end{equation}

\subsection{Weight calibration metric}
In addition to F1 (edge detection), weight recovery is measured by the
Pearson correlation $R$ and the linear regression slope between true and inferred
off-diagonal weights.  A well-calibrated model has slope $\approx 1$; slope $\ll 1$
indicates systematic weight deflation, typically caused by over-aggressive spectral
rescaling.

% ============================================================
\section{Results: $K=2$ (\texttt{daleN100\_bc2f18\_8ea3cb})}
% ============================================================

\subsection{Ground-truth vs.\ recovered connectivity}

\begin{figure}[H]
  \centering
  \includegraphics[width=\textwidth]{figures/k2_connectivity_matrix.jpg}
  \caption{\textbf{Ground-truth vs.\ recovered connectivity matrix ($K=2$).}
    Both panels display the full $100\times100$ synaptic weight matrix $A$, with
    postsynaptic neuron index on the $y$-axis and presynaptic neuron index on the
    $x$-axis.  Warm colours (red) indicate positive (excitatory) weights; cool colours
    (blue) indicate negative (inhibitory) weights; the colour scale is symmetric about
    zero.  A horizontal dashed line at index~80 marks the excitatory/inhibitory
    boundary: rows~0--79 are excitatory neurons (off-diagonal entries $\ge 0$),
    rows~80--99 are inhibitory neurons (off-diagonal entries $\le 0$).
    \textbf{Left:} ground-truth matrix $A_{\mathrm{true}}$, which is sparse and
    Dale-consistent by construction.  The upper band (rows~80--99) shows a dense block
    of strong negative (blue) weights reflecting the inhibitory neurons' outgoing
    connections.  The excitatory block (rows~0--79) is visibly sparser with moderate
    positive weights.  Self-connections on the diagonal are uniformly inhibitory (blue
    stripe).
    \textbf{Right:} recovered matrix $\hat{A}$ after training.  The overall sparsity
    pattern closely matches the ground truth, with the inhibitory block recovered at
    near-perfect F1 and the excitatory block also well recovered.
    Edge-detection statistics: TP$=1338$, FP$=13$, FN$=1$, $F_1=0.995$, $R=0.990$.}
  \label{fig:k2_connectivity}
\end{figure}

\FloatBarrier

\subsection{Connectivity diagnostics}

\begin{figure}[H]
  \centering
  \includegraphics[width=\textwidth]{figures/k2_diagnostics.jpg}
  \caption{\textbf{Connectivity diagnostics ($K=2$).}
    \textbf{Top-left:} bar chart of edge classification counts showing 1338 true
    positives, 13 false positives, and 1 false negative, with inset summary statistics:
    precision$=0.990$, recall$=0.999$, $F_1=0.995$.  Per-block F1 is $0.995$ for
    both excitatory and inhibitory blocks.
    \textbf{Top-centre:} edge-prediction confusion map on the $100\times100$ neuron
    grid.  Green dots are true-positive edges, magenta are false positives, red are
    false negatives, grey are true negatives.  The near-absence of magenta and red
    confirms very few misclassified edges, with no systematic spatial pattern to
    the errors.
    \textbf{Top-right:} off-diagonal weight correlation scatter ($R=0.990$,
    slope$=1.000$).  Each point is one inferred weight plotted against the
    corresponding true weight.  Points cluster tightly along the identity line,
    confirming that the model recovers both edge structure and weight magnitudes
    without systematic bias.
    \textbf{Bottom-left:} log-baseline bias scatter for state~0 ($R=0.999$, blue).
    \textbf{Bottom-right:} log-baseline bias scatter for state~1 ($R=0.999$, red).
    Both scatter plots show inferred $\hat{B}_{k,n}$ tightly concentrated on the
    identity line against true $B_{k,n}^{\mathrm{true}}$, indicating accurate recovery
    of the state-specific baseline excitability for each neuron in both states.}
  \label{fig:k2_diagnostics}
\end{figure}

\FloatBarrier

\subsection{State prediction}

\begin{figure}[H]
  \centering
  \includegraphics[width=\textwidth]{figures/k2_state_prediction.jpg}
  \caption{\textbf{State prediction ($K=2$).}
    \textbf{Top row, left:} mean state occupancy bar chart showing that the two
    hidden states have approximately equal occupancy ($\approx 0.50$ each),
    confirming that neither state collapsed to zero during training.
    \textbf{Top row, centre:} per-state accuracy bar chart; both State~0 and State~1
    are recovered with accuracy $\approx 0.99$.
    \textbf{Top row, right:} confidence histogram showing the distribution of
    $\max_k\gamma_{t,k}$ across all time bins.  The mass is strongly concentrated
    near $1.0$, indicating that the model is highly confident in its state
    assignments at almost every time step; overall accuracy $= 0.992$.
    \textbf{Bottom:} state trace over 1000\,s.  The true hidden state (black)
    alternates between State~0 and State~1 with dwell times of several seconds.
    The inferred state (orange, $\arg\max_k \gamma_{t,k}$) tracks the true sequence
    almost perfectly, with only rare brief disagreements at state boundaries.}
  \label{fig:k2_state}
\end{figure}

\FloatBarrier

\subsection{Training monitor}

\begin{figure}[H]
  \centering
  \includegraphics[width=\textwidth]{figures/k2_training_monitor.png}
  \caption{\textbf{Training monitor ($K=2$).}
    \textbf{Panel 1 -- E-step NLL:} the E-step negative log-likelihood per bin drops
    steeply from $\approx 3.5\times10^7$ at EM iteration~0 to $\approx 1.0\times10^7$
    by iteration~3, then flattens and converges smoothly over the remaining 22
    iterations.  The rapid initial drop reflects the warm-start initialisation;
    the flat tail confirms convergence.
    \textbf{Panel 2 -- Spectral radius:} $\rho(A)$ begins at $\approx 4.0$ after
    the warm start, then falls sharply within the first $\sim$100 M-step epochs as
    spectral rescaling is applied every 3 epochs.  After epoch $\sim$100, $\rho(A)$
    stabilises just below the dashed $\rho_{\max}=0.92$ threshold and remains there
    for the rest of training, confirming that the spectral stability constraint is
    enforced throughout.
    \textbf{Panel 3 -- Connectivity support:} the number of nonzero off-diagonal
    edges starts high ($\approx 7500$) from the dense warm-start, then decreases
    rapidly as the proximal operator prunes weak edges.  A pronounced second drop
    around epoch~500 corresponds to the cadence switch from $q_{\mathrm{early}}=8$
    to $q_{\mathrm{late}}=2$, after which finer-grained pruning produces the final
    sparse support of $\approx 4500$ edges.
    \textbf{Panel 4 -- M-step NLL and L1 norm:} the Poisson NLL (blue, left axis)
    and the L1 norm of off-diagonal weights (red dashed, right axis) both decrease and
    converge over the 750 M-step epochs, confirming simultaneous data fitting and
    sparsity enforcement.}
  \label{fig:k2_training}
\end{figure}

\FloatBarrier

\subsection{Weight and bias residuals}

\begin{figure}[H]
  \centering
  \includegraphics[width=\textwidth]{figures/k2_residuals.jpg}
  \caption{\textbf{Edge and weight residuals ($K=2$).}
    \textbf{Panel 1 -- $A_{\mathrm{off}}$ inhibitory residual:} histogram of
    $\hat{A}_{ij} - A_{ij}^{\mathrm{true}}$ for inhibitory off-diagonal entries.
    Centred near zero (median~$\approx -0.005$, std~$\approx 0.025$), with no heavy
    tails, confirming accurate inhibitory weight recovery.
    \textbf{Panel 2 -- $A_{\mathrm{off}}$ excitatory residual:} same for excitatory
    off-diagonal entries.  Similarly centred (median~$\approx 0.002$,
    std~$\approx 0.016$) and symmetric, confirming no systematic excitatory weight
    deflation from the L2 regulariser.
    \textbf{Panel 3 -- $A_{\mathrm{diag}}$ residual:} residuals for diagonal
    (self-)weights, approximately symmetric around zero (median~$\approx 0.003$),
    confirming accurate self-inhibition recovery.
    \textbf{Panel 4 -- Edge accuracy bar chart:} counts of TP (1338), FP (13), and
    FN (1) graphically confirm near-perfect edge classification.
    \textbf{Panel 5 -- $B$ low-rate residual:} residuals for neurons with low
    baseline rates ($B < 2.95$), tightly centred near zero.
    \textbf{Panel 6 -- $B$ high-rate residual:} residuals for neurons with high
    baseline rates ($B \ge 2.95$), also centred near zero, confirming that bias
    recovery is accurate across the full dynamic range of baseline firing rates.}
  \label{fig:k2_residuals}
\end{figure}

\FloatBarrier

% ============================================================
\section{Results: $K=3$ (\texttt{daleN100\_f7d754\_ce2202})}
% ============================================================

\subsection{Ground-truth vs.\ recovered connectivity}

\begin{figure}[H]
  \centering
  \includegraphics[width=\textwidth]{figures/k3_connectivity_matrix.jpg}
  \caption{\textbf{Ground-truth vs.\ recovered connectivity matrix ($K=3$).}
    Same layout as Figure~\ref{fig:k2_connectivity} for dataset
    \texttt{daleN100\_f7d754\_ce2202}.  The ground-truth excitatory block
    (rows~0--79) is sparser than in the $K=2$ dataset, with fewer strong positive
    weights; the inhibitory block (rows~80--99) is again densely connected with strong
    negative weights.  The recovered matrix (right) reproduces both support and weight
    magnitudes with high fidelity.  Zero false negatives (FN$=0$) means every true
    edge is detected; the nine false positives represent spurious edges of small
    magnitude.  Edge-detection statistics: TP$=1358$, FP$=9$, FN$=0$,
    $F_1=0.997$, $R=0.988$.  The sparser excitatory ground truth compared to the
    $K=2$ case explains why excitatory F1 reaches $0.997$ here: fewer indirect
    two-hop paths create confounding lag-1 correlations.}
  \label{fig:k3_connectivity}
\end{figure}

\FloatBarrier

\subsection{Connectivity diagnostics}

\begin{figure}[H]
  \centering
  \includegraphics[width=\textwidth]{figures/k3_diagnostics.jpg}
  \caption{\textbf{Connectivity diagnostics ($K=3$).}
    Same layout as Figure~\ref{fig:k2_diagnostics}.
    \textbf{Top-left:} TP$=1358$; FP$=9$ and FN$=0$ are barely visible,
    confirming near-perfect edge recovery.  Summary: precision$=0.993$,
    recall$=1.000$, $F_1=0.997$.
    \textbf{Top-centre:} confusion map showing an almost entirely green field of
    true positives, with only a handful of magenta false-positive dots and no red
    false-negative dots, reflecting the zero-FN result.
    \textbf{Top-right:} weight correlation scatter ($R=0.988$, slope$=0.956$).
    Points cluster tightly along a line of slope $0.956$, slightly below the
    identity, indicating mild weight deflation consistent with the spectral
    rescaler slightly undershooting true magnitudes.  This is the primary source
    of imperfection in the $K=3$ run.
    \textbf{Bottom row:} log-baseline bias scatter for all three states
    (State~0 blue, State~1 red, State~2 green), each with $R=0.999$.  All three
    scatter plots show tight concentration on the identity line, confirming that the
    model correctly identifies the distinct baseline firing rate of each neuron across
    all three brain states.}
  \label{fig:k3_diagnostics}
\end{figure}

\FloatBarrier

\subsection{State prediction}

\begin{figure}[H]
  \centering
  \includegraphics[width=\textwidth]{figures/k3_state_prediction.jpg}
  \caption{\textbf{State prediction ($K=3$).}
    Same layout as Figure~\ref{fig:k2_state} extended to three states.
    \textbf{Top row, left:} mean occupancy bar chart showing approximately equal
    occupancy across all three states ($\approx 0.33$ each), confirming balanced
    state usage with no collapsed states.
    \textbf{Top row, centre:} per-state accuracy bar chart: State~0 $=0.99$,
    State~1 $=0.98$, State~2 $=0.99$; overall accuracy $= 0.988$.
    \textbf{Top row, right:} confidence histogram showing $\max_k\gamma_{t,k}$
    strongly concentrated near $1.0$, indicating high-confidence state assignments
    throughout.
    \textbf{Bottom:} state trace over 1000\,s.  The true state sequence (black)
    cycles among all three states with comparable dwell times.  The inferred sequence
    (orange) tracks it closely, with occasional brief errors at state boundaries.
    The three-level structure is clearly discriminated, demonstrating that the model
    scales successfully to $K=3$ without state confusion or collapse.}
  \label{fig:k3_state}
\end{figure}

\FloatBarrier

\subsection{Training monitor}

\begin{figure}[H]
  \centering
  \includegraphics[width=\textwidth]{figures/k3_training_monitor.png}
  \caption{\textbf{Training monitor ($K=3$).}
    Same four-panel layout as Figure~\ref{fig:k2_training}.
    \textbf{Panel 1 -- E-step NLL:} drops from $\approx 3.0\times10^7$ at
    iteration~0 to $\approx 1.0\times10^7$ within the first 3 iterations, then
    converges smoothly, mirroring the $K=2$ behaviour.
    \textbf{Panel 2 -- Spectral radius:} $\rho(A)$ begins at $\approx 3.5$, falls
    to below $\rho_{\max}=0.92$ by epoch $\sim$100, and remains stable thereafter.
    The convergence trajectory is essentially identical to the $K=2$ case, confirming
    that the spectral constraint operates the same way regardless of the number of
    states.
    \textbf{Panel 3 -- Connectivity support:} nonzero edge count decreases from
    $\approx 7500$ to $\approx 4500$ with the same two-phase pattern as $K=2$,
    as the underlying network architecture and sparsity level are the same across
    datasets.
    \textbf{Panel 4 -- M-step NLL and L1 norm:} both quantities decrease and
    stabilise over 750 M-step epochs.  The L1 norm stabilises at a slightly higher
    value than in the $K=2$ run, consistent with the $K=3$ dataset having more
    true edges in total across all blocks.}
  \label{fig:k3_training}
\end{figure}

\FloatBarrier

\subsection{Weight and bias residuals}

\begin{figure}[H]
  \centering
  \includegraphics[width=\textwidth]{figures/k3_residuals.jpg}
  \caption{\textbf{Edge and weight residuals ($K=3$).}
    Same six-panel layout as Figure~\ref{fig:k2_residuals}.
    \textbf{Panel 1 -- $A_{\mathrm{off}}$ inhibitory residual:} centred near zero
    (median~$\approx -0.016$, std~$\approx 0.019$), comparable to the $K=2$ result.
    \textbf{Panel 2 -- $A_{\mathrm{off}}$ excitatory residual:} distribution slightly
    shifted negative (median~$\approx -0.003$, std~$\approx 0.014$), reflecting the
    mild slope$<1$ weight deflation also visible in the weight scatter.
    \textbf{Panel 3 -- $A_{\mathrm{diag}}$ residual:} centred near zero
    (median~$\approx 0.004$), confirming accurate diagonal recovery.
    \textbf{Panel 4 -- Edge accuracy bar chart:} TP$=1358$, FP$=9$, FN$=0$ ---
    the perfect-recall result is clearly visible as the absence of any FN bar.
    \textbf{Panels 5--6 -- $B$ residuals:} low-rate and high-rate bias residuals
    are both tightly distributed around zero (std $<0.04$), confirming accurate
    bias recovery across all three states and the full range of baseline firing rates.}
  \label{fig:k3_residuals}
\end{figure}

\FloatBarrier

% ============================================================
\section{Hyperparameter Reference}
% ============================================================

\small
\begin{longtable}{>{\raggedright\arraybackslash}p{0.20\textwidth}
                  >{\raggedright\arraybackslash}p{0.14\textwidth}
                  >{\raggedright\arraybackslash}p{0.58\textwidth}}
\toprule
Name & Value & Role \\
\midrule
\endhead

\multicolumn{3}{l}{\textit{Sparsity (swept)}} \\[2pt]
$\lambda_1$ (\texttt{l1}) & $10^{-3}$ (best) & Base L1 strength; scales the initial proximal thresholds $\tau_0$ for both blocks. \\
$\lambda_2$ (\texttt{l2}) & $10^{-4}$ (best) & L2 weight decay; keeps weights bounded and aids numerical stability. \\
\midrule

\multicolumn{3}{l}{\textit{EM schedule}} \\[2pt]
\texttt{em\_iters} & 25 & Number of EM outer iterations. \\
\texttt{m\_epochs} & 30 & Adam epochs per EM round. \\
\midrule

\multicolumn{3}{l}{\textit{Optimiser}} \\[2pt]
\texttt{lr} & $5\times10^{-3}$ & Adam learning rate. \\
\texttt{batch\_size} & 2048 & Minibatch size for stochastic gradient estimation. \\
\texttt{grad\_clip} & 1.0 & Gradient norm clipping threshold. \\
\midrule

\multicolumn{3}{l}{\textit{Proximal / adaptive $\tau$ controller}} \\[2pt]
\texttt{prox\_scale} $(\alpha_{\mathrm{prox}})$ & $5\times10^{-3}$ & Multiplier converting $\lambda_1$ to the initial proximal threshold $\tau_0$. \\
\texttt{tau\_adapt\_start\_frac} & 0.25 & Fraction of total Phase~1 epochs before tau adaptation begins (warm-up). \\
\texttt{tau\_gain} $(g_\tau)$ & 0.05 & Multiplicative gain for FP/FN feedback in tau update. \\
\texttt{tau\_min\_mult} & 0.2 & Lower clamp: $\tau_{\min}=0.2\,\tau_0$. \\
\texttt{tau\_max\_mult} & 6.0 & Upper clamp: $\tau_{\max}=6.0\,\tau_0$. \\
\texttt{proj\_every\_early} & 8 & Proximal+Dale projection applied every $n$ epochs (early training). \\
\texttt{proj\_every\_late} & 2 & Proximal+Dale projection applied every $n$ epochs (late training). \\
\texttt{proj\_switch\_frac} & 0.65 & Fraction of Phase~1 epochs at which cadence switches from early to late. \\
\midrule

\multicolumn{3}{l}{\textit{Structural constraints}} \\[2pt]
\texttt{diag\_inhibitory\_floor} $(d_{\min})$ & 0.05 & Minimum magnitude of self-inhibition on the diagonal. \\
spectral $\rho_{\max}$ & 0.92 & Spectral radius threshold above which $A$ is rescaled. \\
spectral $\rho_{\mathrm{target}}$ (training) & 0.90 & Target spectral radius during training rescaling. \\
spectral cadence & every 3 epochs & How often spectral rescaling is applied during training. \\
spectral $\rho_{\mathrm{target}}$ (final) & 0.919 & Target for the single post-training spectral stabilisation step. \\
\midrule

\multicolumn{3}{l}{\textit{Warm start}} \\[2pt]
\texttt{warm\_ridge} $(\lambda_{\mathrm{ridge}})$ & $10^{-2}$ & Ridge regulariser in the warm-start linear solve. \\
\texttt{warm\_keep\_ratio} $(r_{\mathrm{keep}})$ & 0.15 & Fraction of off-diagonal entries retained after warm-start thresholding. \\
\texttt{b\_init\_noise} $(\sigma_B)$ & 0.05 & Std of Gaussian noise added to bias warm start to break state symmetry. \\
\midrule

\multicolumn{3}{l}{\textit{HMM priors}} \\[2pt]
\texttt{pi\_pseudocount} $(\alpha_\pi)$ & $10^{-2}$ & Additive smoothing on initial state distribution. \\
\texttt{P\_pseudocount} $(\epsilon_P)$ & $10^{-3}$ & Additive smoothing on transition counts. \\
\texttt{delta\_t} $(\Delta t)$ & 0.01 & Time bin width in seconds. \\
\midrule

\multicolumn{3}{l}{\textit{Architectural}} \\[2pt]
\texttt{n\_lags} $(L)$ & 1 & Number of lagged time steps in the design matrix. \\
$N$ & 100 & Total number of neurons. \\
$K$ & 2 or 3 & Number of hidden states (dataset-dependent). \\
$\mathrm{exc\_ratio}$ & 0.8 & Fraction of neurons treated as excitatory; $n_{\mathrm{exc}}=\lfloor 0.8N\rfloor = 80$. \\
$T$ & 200{,}000 & Time bins per dataset. \\
\bottomrule
\end{longtable}
\normalsize

\subsection*{Sweep ranges tried}

\begin{tabular}{ll}
\toprule
Parameter & Values swept \\
\midrule
$\lambda_1$ & $\{10^{-3},3\times10^{-3},10^{-2},3\times10^{-2},10^{-1}\}$ \\
$\lambda_2$ & $\{10^{-4},3\times10^{-4},10^{-3},3\times10^{-3},10^{-2}\}$ \\
\texttt{n\_lags} & $\{1,2\}$ ($L=2$ tested; reverted: FP dropped but FN rose, net F1 unchanged) \\
\texttt{tau\_gain} & $\{0.05, 0.06, 0.10\}$ (0.10 caused oscillation; 0.05 in use) \\
\texttt{tau\_max\_mult} & $\{5.0, 6.0, 10.0\}$ (10.0 caused limit cycle; 6.0 in use) \\
\texttt{proj\_every\_late} & $\{1, 2, 3\}$ (1 caused oscillation; 2 in use) \\
spectral cadence & \{every epoch, every 3 epochs\} (every epoch caused weight deflation) \\
spectral target & $\{0.919, 0.90\}$ (0.919 gave insufficient weight-growth buffer) \\
\bottomrule
\end{tabular}

% ============================================================
\section{Known Limitations}
% ============================================================

\paragraph{Excitatory block ceiling ($K=2$).}
With lag-1 features, indirect excitatory paths $i\!\to\!k\!\to\!j$ produce
same-sign, similar-magnitude lag-1 correlations as direct connections $i\!\to\!j$.
Dale's law eliminates inhibitory confounders but not excitatory ones.  All post-hoc
correction methods tried (network deconvolution, two-hop subtraction, path penalty)
were unable to break this degeneracy.  The F1 ceiling for the excitatory block
appears to be $\approx0.80$ on datasets with denser excitatory ground truth.
The $K=3$ dataset has a sparser excitatory ground truth (fewer true edges, hence fewer
indirect two-hop paths), which explains the substantially higher excitatory F1 of $0.997$.

\paragraph{Weight calibration ($B$).}
The log-baseline bias vectors $B_k$ are recovered with $R \ge 0.999$ per state for
both $K=2$ and $K=3$.  However, the scatter plots show a systematic offset: inferred
values lie slightly below the true values, consistent with the L2 regulariser
shrinking the biases toward zero.  This matters for control applications where the
absolute baseline rate of each neuron is needed.

\end{document}
